\documentclass{article}
\usepackage[T1]{fontenc}
\usepackage[margin=1in]{geometry}
\usepackage{graphicx}
\usepackage{amsmath}
\usepackage{amssymb}
\usepackage{booktabs}
\usepackage{tabularx}
\usepackage{xurl}
\usepackage[hidelinks]{hyperref}

\newcommand{\zscore}{\operatorname{zscore}}
\newcommand{\sigmoid}{\operatorname{sigmoid}}
\newcolumntype{L}{>{\raggedright\arraybackslash}X}
\setlength{\emergencystretch}{2em}

% Add the named PDF/PNG figures when they are ready. Until then this macro
% displays a labelled placeholder, so the document compiles without images.
% For SVG artwork, export to PDF or replace the call with \includesvg.
\newcommand{\projectfigure}[2]{%
  \IfFileExists{#1}{%
    \includegraphics[width=0.92\linewidth]{#1}%
  }{%
    \fbox{\parbox[c][4cm][c]{0.86\linewidth}{%
      \centering\textbf{Figure placeholder}\par\medskip
      #2\par\medskip\small\path{#1}%
    }}%
  }%
}

\title{CorrPFN: An EPIT-Informed Prior for Pitting-Potential Prediction}
\author{}
\date{}

\begin{document}
\maketitle

\section{Introduction}
\label{sec:introduction}

Pitting corrosion is a localized degradation process in which a protective
passive film breaks down and stable pits begin to form. A commonly reported
measure of this transition is the pitting potential, denoted here by EPIT.
Under comparable environmental and experimental conditions, a more positive
EPIT generally indicates that a higher applied potential is required to
initiate sustained pitting and therefore corresponds to greater resistance to
passive-film breakdown. Predicting EPIT from alloy composition, exposure
conditions, and test information is nevertheless difficult because the
available data are heterogeneous, sparse relative to the number of possible
material--environment combinations, and strongly affected by measurement
protocol.

This work investigates whether a Prior-Data Fitted Network can be specialized
for EPIT prediction by modifying the synthetic task distribution used during
pretraining. The starting point is TabICL, which learns an in-context
prediction procedure from many synthetically generated tabular tasks. Its
generic prior produces diverse nonlinear and tree-structured prediction
problems, but its features are scientifically anonymous and its target has no
fixed physical interpretation. It therefore does not explicitly represent the
material, environment, and process structure of an EPIT dataset or the
physical direction of a pitting-potential target.

CorrPFN introduces an informed synthetic prior that combines empirical
physical features with stochastic prediction tasks. Alloy compositions are
sampled from observed Fe and Ni--Cr--Mo templates, perturbed, and combined with
empirical environmental conditions. An MLP or tree SCM maps the resulting
features to a synthetic target. Independently, one of seven calibrated EPIT
rules supplies a corrosion-specific target component. Both components are
standardized and combined using an optimized mixture weight. Generic tasks
remain available in the same pretraining run.

The added component represents broad probabilistic tendencies rather than a
deterministic corrosion model: material passivity tends to increase EPIT,
environmental aggressiveness tends to decrease it, material--environment
interactions can further reduce it, and test protocol can shift the reported
value in either direction. Empirical feature generation constrains the inputs
to patterns represented in the source dataset, while the SCM component
retains nonlinear task diversity beyond the analytical rules.

The final workflow separates target-rule calibration and model selection on
608 development rows from evaluation on 152 rows assigned to different
composition groups. The selected CorrPFN checkpoint obtains a final-test
Spearman correlation of 0.7681 and MAE of 186.8~mV, improving on the reported
pretrained TabICL V2, CatBoost, and local generic-prior references. These are
internal results from one source dataset and one final training seed; the
scope of dataset exposure and the limits of the comparison are stated below.

This report describes the final implementation, identified as v7 in the
experiment artifacts. Earlier experiments are discussed only where they
explain a retained design choice or affect the interpretation of the results.

\section{EPIT Dataset and Evaluation Design}
\label{sec:epit-dataset}

\subsection{Source Data, Rows, and Feature Schema}

The evaluated EPIT task is derived from the public dataset
\textit{Electrochemical Metrics for Corrosion Resistant Alloys}, released on
Figshare and described by Nyby et al. in \textit{Scientific Data}
\cite{nyby2021electrochemical}. In this project, the dataset is identified as
\path{electrochemical_metrics_alloys}, and the evaluated task identifier is
\path{electrochemical_metrics_alloys__pitting_potential__epit_mv_sce_avg}.
The source workbook table is \texttt{Pitting Potential}, and the prediction
target is \texttt{Epit, mV (SCE) Avg.}, where SCE denotes the saturated calomel
reference electrode.

The source table contains \(810\) observed rows. Of these, \(760\) have a
usable numeric value for the selected average EPIT target. The remaining
\(50\) rows are excluded because the target is missing or cannot be
interpreted as a numeric EPIT measurement.

The evaluated input schema contains \(21\) features:
\begin{equation}
21\text{ features}
=17\text{ material}+3\text{ environment}+1\text{ process/history}.
\label{eq:epit-schema}
\end{equation}
This subset is aligned with the feature set used by the reference modelling
study, allowing results to be compared under a consistent input definition.
Columns omitted from the source workbook should not therefore be interpreted
as scientifically irrelevant.

The \(17\) material features are alloy-composition values reported in weight
percent, retained in the following fixed order:
\[
\mathrm{Fe,\ Cr,\ Ni,\ Mo,\ W,\ Nb,\ Al,\ V,\ Ta,\ Re,\ Ce,\ Ti,\ Co,\ B,\ Mg,\ Y,\ Gd}.
\]
They are retained as separate inputs rather than being reduced to one
aggregate alloy descriptor.

The next three columns are \texttt{Test Temp. oC}, the test temperature in
degrees Celsius; \texttt{[Cl-] M}, the chloride-ion concentration in molar
units; and \texttt{[Cl-] pH}, the reported pH of the chloride-containing
solution. The final column is \texttt{[Cl-] Test Method}, which records the
experimental method or protocol associated with the measurement.

Test method is classified as process/history information because polarization
procedure, scan conditions, specimen preparation, stabilization time, and the
operational breakdown criterion can all affect the reported EPIT. It is a
single recorded protocol feature; these individual experimental factors are
not all separately observed in the retained schema.

\subsection{Target Meaning and Experimental Heterogeneity}

EPIT is the reported electrochemical potential associated with passive-film
breakdown or the onset of sustained pitting. Under comparable
reference-electrode, environmental, surface, and test-protocol conditions, a
more positive EPIT indicates that a higher applied potential is required
before the reported breakdown criterion is reached. It is therefore generally
interpreted as stronger resistance to pit initiation.

This interpretation is conditional rather than universal. The workbook
combines measurements obtained using potentiodynamic, potentiostatic,
scratch-based, and other polarization procedures. Different studies can
identify breakdown using different current thresholds, scan rates,
stabilization periods, or visual criteria. Consequently, the operational
meaning of the reported EPIT is not identical for every row.

A measured EPIT depends jointly on alloy composition and microstructure,
chloride concentration and other solution chemistry, pH and temperature,
surface preparation and prior exposure, polarization method and scan rate,
and the criterion used to identify breakdown. A higher raw EPIT therefore
does not by itself prove that one alloy is intrinsically more resistant than
another when the measurements were obtained under different environments or
protocols.

\subsection{Data Representation}
\label{sec:data-representation}

The \(20\) material and environment columns are converted to numeric values.
During real-data model evaluation, numeric mean imputation and categorical
encoding are fitted on the context/training portion and applied to the
validation or final-test portion. The local preprocessing implementation
restores the original column order after these transformations.

Test method is represented by a single integer-coded column rather than
one-hot encoding. The complete task contains \(52\) categories including a
missing-value category; the empirical synthetic profile stores a fixed
lexicographically ordered mapping over these categories. Context-fitted
evaluation encodings can contain fewer categories, with an explicit code for
unknown or missing values. These integers are nominal identifiers. They do
not define a physical ordering or distance between experimental methods.
Treating the encoded column as numeric remains an approximation.

The synthetic feature profile has its own stored imputation means and category
mapping, described in Section~\ref{sec:feature-profile-provenance}. These
profile statistics must be distinguished from context-fitted preprocessing
during model evaluation. Similarly, analytical target-rule calibration fits
its imputation and scaling state within the relevant calibration context.

\subsection{Composition-Aware Split Construction}
\label{sec:composition-split}

The split is designed primarily to evaluate prediction for unseen alloy
composition groups. Measurements of sufficiently similar compositions are
therefore kept together. Grouping uses the \(17\) composition columns visible
to TabICL. The compositions are not renormalized to \(100\%\), because seven
composition columns from the source data are absent from the fixed model
schema. For grouping only, missing values are treated as zero and each visible
composition value is rounded to \(0.01\)~wt.\%.

For two rounded composition vectors \(x_i\) and \(x_j\), their distance is
\[
d(i,j)=\sum_{e=1}^{17}|x_{i,e}-x_{j,e}|.
\]
Every pair satisfying
\[
d(i,j)\leq1.0\ \text{wt.\%}
\]
is connected. The final groups are the connected components of this threshold
graph. Consequently, every directly similar pair remains in the same component
and cannot cross either the development/final-test boundary or a development
fold boundary.

This connected-component definition is intentionally transitive. For example,
if alloy \(A\) is within \(1\)~wt.\% of \(B\), and \(B\) is within
\(1\)~wt.\% of \(C\), all three are grouped together even when \(A\) and
\(C\) are farther apart. The frozen split contains \(301\) composition
components. The largest component contains \(45\) rows, and the largest
within-component span is \(5.0\)~wt.\%. This wider span is the cost of
guaranteeing that no directly close pair crosses a split boundary.

Complete composition components were assigned to obtain exactly
\[
|\mathcal{D}_{\mathrm{dev}}|=608,
\qquad
|\mathcal{D}_{\mathrm{final}}|=152.
\]
The outer split and the five development folds were balanced using material
class, EPIT range, temperature, chloride concentration, pH, test method,
source reference, and row count. EPIT values were therefore used once to
construct a balanced fixed split; the split is not target-blind. Source
reference overlap is balanced but not prohibited.

\subsection{Development and Final Evaluation}

The development set is divided into five fixed folds. For each fold,
\(486\)--\(487\) rows are used as context or training rows and
\(121\)--\(122\) rows are used as validation rows. Every development row
serves as validation data once.

Three different evaluations must be distinguished:
\begin{itemize}
    \item \textbf{Direct target-rule evaluation} measures how well a corrosion
    formula orders real EPIT values on development validation rows, using
    coefficients fitted on the corresponding context rows.
    \item \textbf{Transformer development evaluation} measures how well a
    synthetically pretrained TabICL checkpoint predicts the five development
    validation folds. It is used for Optuna and checkpoint selection.
    \item \textbf{Final model evaluation} uses all \(608\) development rows
    as context and the \(152\) reserved rows as queries for the locked model.
\end{itemize}

The seven candidate target rules are calibrated on the \(452\) development
rows belonging to the eligible Fe and Ni--Cr--Mo classes. Transformer
development evaluation uses all \(608\) development rows, and final evaluation
uses all \(152\) final-test rows. The restriction of the informed generator
and rule calibration to two alloy families does not restrict the evaluated
task to those families.

Final-test targets are excluded from rule calibration, Optuna selection, and
checkpoint selection in this workflow. The complete corpus had nevertheless
influenced earlier exploratory development, and feature information from it
is retained in the empirical profiles. The resulting evaluation is an
internal composition-separated assessment, with the limitations discussed in
Section~\ref{sec:discussion}.

\begin{figure}[htbp]
    \centering
    \projectfigure{images/corrpfn_split_overview.pdf}
    {Composition groups, the 608/152 partition, and five development folds.}
    \caption{Composition-aware evaluation design. Target-rule calibration and
    model selection use the development partition; the selected model uses
    all development rows as context for the final evaluation.}
    \label{fig:split-overview}
\end{figure}

\section{TabICL Model and Regression Setup}
\label{sec:tabicl}

\subsection{Generic Synthetic Tasks}

TabICL is pretrained on a distribution of synthetic supervised-learning tasks.
In this context, the \emph{prior} is the stochastic procedure that generates
those tasks, not one fixed synthetic dataset. Each sampled task contains a
feature table \(X\), a target \(y\), and a context/query split. During
pretraining, the model receives all feature rows and the targets associated
with the context rows, and it learns to predict the withheld query targets.

The generic TabICL prior is domain-independent. Its purpose is to expose the
model to diverse tabular relationships from which it can learn a reusable
in-context prediction procedure. The principal generator choices are
\texttt{mlp\_scm}, which produces nonlinear relationships using a randomly
sampled multilayer perceptron; \texttt{tree\_scm}, which produces split-based
and piecewise relationships using a random tree or ensemble; and
\texttt{mix\_scm}, which samples between the two generators.

The MLP generator supports two internal topologies through the sampled
parameter \texttt{is\_causal}. When it is true, latent root variables are
transformed into intermediate variables from which the observed features and
target are selected. When it is false, the sampled variables are used directly
as features and the final MLP output becomes the target. The tree generator
uses the latter, direct predictive structure. Conceptually,
\begin{equation}
\begin{aligned}
\text{causal-style MLP:}\quad &z\longrightarrow h\longrightarrow(X,y),\\
\text{direct MLP or tree:}\quad &X=z,\qquad y=f(X).
\end{aligned}
\label{eq:generic-generator-modes}
\end{equation}
Here, \(z\) denotes sampled root variables and \(h\) denotes generated
intermediate variables. The SCM terminology describes the synthetic
construction; it does not identify causal effects in the real corrosion data.

Generic tasks pass through the standard TabICL preprocessing pipeline.
Depending on configuration, this can convert numerical columns into
categorical-like columns, limit extreme values, standardize and permute
features, remove unusable columns, and pad unused positions. Invalid tasks are
rejected and resampled. The informed branch described in
Section~\ref{sec:informed-prior} supplies its own fixed-schema preprocessing.

These mechanisms allow the generic prior to produce nonlinear effects,
interactions, correlated or redundant features, irrelevant variables,
different feature counts, and different context/query sizes. However, its
columns remain scientifically anonymous: no feature is inherently identified
as an alloy-composition variable, an environmental condition, or a test
parameter.

\subsection{PFN Interpretation}

Let \(\theta\) denote the complete latent mechanism sampled for one synthetic
task. It can include the generator family, network or tree parameters, feature
distributions, dependency structure, noise, and other task-specific random
quantities. A synthetic dataset is then sampled conditionally on that
mechanism. Thus, \(\theta\) describes how one task is generated; it is not the
generated dataset itself.

The TabICL transformer has a separate set of shared parameters, denoted by
\(\phi\). During pretraining, many mechanisms \(\theta\) and their associated
datasets are sampled, while \(\phi\) is optimized across all tasks.

For a labeled context dataset
\[
\mathcal{D}_{\mathrm{c}}=\{(x_i,y_i)\}_{i=1}^{n}
\]
and a query input \(x^\ast\), the posterior predictive distribution induced
by a task prior \(p(\theta)\) is
\begin{equation}
p(y^\ast\mid x^\ast,\mathcal{D}_{\mathrm{c}})
=\int p(y^\ast\mid x^\ast,\theta)
       p(\theta\mid\mathcal{D}_{\mathrm{c}})\,d\theta.
\label{eq:posterior-predictive}
\end{equation}
A PFN is trained to amortize this inference. Its predictive distribution is
intended to approximate the posterior predictive distribution associated with
its synthetic training prior:
\begin{equation}
q_{\phi}(y^\ast\mid x^\ast,\mathcal{D}_{\mathrm{c}})
\approx p(y^\ast\mid x^\ast,\mathcal{D}_{\mathrm{c}}).
\label{eq:pfn-approximation}
\end{equation}

At inference time, the transformer produces predictions through forward
evaluation. Its parameters \(\phi\) remain fixed: it does not retrain on the
context data, enumerate possible mechanisms, or numerically evaluate the
integral in Equation~\eqref{eq:posterior-predictive}. The reported inference
configuration uses an ensemble of eight estimators.

This interpretation is conditional on the synthetic training distribution.
The model is not guaranteed to approximate the unknown posterior predictive
distribution of real corrosion data. Transfer depends on how well the
synthetic prior covers the statistical structure of the real task, as well as
on model capacity, training diversity, optimization, and the amount and
quality of context data.

\subsection{Regression and Training Scope}

The local models use the available TabICL V1 implementation with a
continuous-target regression path added for this project. This path is not
native to the original TabICL V1 implementation. Synthetic targets are
standardized for regression rather than converted into class labels. At
evaluation, predictions are returned on the EPIT scale and the reported point
output is the predictive median.

All locally trained models considered here use Stage~1 of the original TabICL
curriculum. Stage~1 trains from scratch on synthetic tables containing up to
\(1{,}024\) rows. Later curriculum stages extend long-context capability and
were unnecessary for the \(760\)-row EPIT task. The proxy and final training
horizons are specified in Section~\ref{sec:model-selection}.

\section{CorrPFN Synthetic Prior}
\label{sec:informed-prior}

\subsection{Generation Pipeline}

The informed prior separates the generation of physical features from the
construction of a synthetic response. Each informed task first receives a
complete empirical physical table
\[
X_{\mathrm{raw}}\in\mathbb{R}^{n\times21}.
\]
Its standardized representation, \(X_{\mathrm{std}}\), is used by both the
SCM target head and the transformer. The raw representation is retained for
the analytical EPIT rule:
\begin{equation}
X_{\mathrm{raw}}\longrightarrow
\begin{cases}
g_k(X_{\mathrm{raw}}), & \text{analytical EPIT rule},\\
f_{\theta}(X_{\mathrm{std}}), & \text{SCM target head},\\
X_{\mathrm{std}}, & \text{transformer input}.
\end{cases}
\label{eq:informed-pipeline}
\end{equation}
The two target components are then mixed and standardized. This order ensures
that the SCM response is generated from the same physical columns that the
transformer observes, while the analytical rule can use their original units.

\begin{figure}[htbp]
    \centering
    \projectfigure{images/corrpfn_generation_pipeline.pdf}
    {Empirical features, raw and standardized representations, SCM and EPIT
    target branches, target mixture, and generic/informed routing.}
    \caption{Final CorrPFN synthetic-task construction. The empirical feature
    generator provides the inputs, the SCM contributes stochastic nonlinear
    responses, and the calibrated EPIT rule contributes domain structure.}
    \label{fig:generation-pipeline}
\end{figure}

\subsection{Empirical Alloy-Composition Generator}
\label{sec:composition-generator}

Compositions are sampled from real templates in the versioned
\texttt{epit\_dataset\_v1} profile. The complete profile contains \(403\)
distinct reported composition vectors. The final informed generator uses the
restricted bank of \(298\) Fe-alloy templates and \(17\) Ni--Cr--Mo templates.
Family probabilities are proportional to these distinct-template counts:
\[
P(\mathrm{Fe})=\frac{298}{315}\approx0.9460,
\qquad
P(\mathrm{Ni\!-!Cr\!-!Mo})=\frac{17}{315}\approx0.0540.
\]
These are frequencies among distinct templates, not among all experimental
rows. A template is sampled within the selected family. The runtime generator
loads the static profile and does not reopen or reprocess the source Excel
workbook during training.

A selected template is represented using the complete \(24\)-element source
composition. For every positive component \(c_{ij}>0\), a multiplicative
perturbation is applied:
\begin{equation}
\widetilde c_{ij}=c_{ij}\exp(\epsilon_{ij}),
\qquad \epsilon_{ij}\sim\mathcal{N}(0,\tau^2),
\label{eq:composition-perturbation}
\end{equation}
where \(\tau\) is the composition perturbation strength. Zero and unreported
elements remain zero. The perturbation therefore changes the amounts of
reported alloying elements without introducing elements absent from that
template.

The complete \(24\)-element composition is subsequently closed to \(100\)
weight percent:
\begin{equation}
c_{ij}^{\mathrm{closed}}
=100\frac{\widetilde c_{ij}}{\sum_{e=1}^{24}\widetilde c_{ie}}.
\label{eq:composition-closure}
\end{equation}
Only after this closure are the \(17\) elements visible to the model selected.
Consequently, the sum of the visible elements can be below \(100\%\) when
part of the source composition belongs to elements omitted from the evaluated
schema. At \(\tau=0\), the generator replays a template after this closure;
it does not necessarily preserve an unnormalized source total.

This approach retains dominant elements, approximate concentration ranges,
sparsity, and element co-occurrence patterns represented in the template bank.
The joint material structure is supplied by these templates and full-vector
closure. Additional shared block coupling and masked Dirichlet generation are
disabled in this branch.

\subsection{Empirical Environmental Conditions}

The environmental part of each synthetic row is generated separately from its
composition. One of the \(565\) real Fe/Ni--Cr--Mo experimental records in the
versioned \texttt{epit\_fe\_nicrmo\_features\_v1} profile is sampled with
replacement. Temperature, chloride concentration, pH, and test method are
retained as a complete tuple from the same source row. This preserves
experimentally observed combinations of environmental conditions.

Composition and environment are sampled independently. The generator therefore
retains empirical composition structure and empirical environmental structure
without copying the dataset's existing association between a particular alloy
family and a particular experimental campaign. The resulting composition and
environment combination need not itself occur in the source workbook.

Missing environmental values use the means stored in the feature profile.
Test methods use its fixed \(52\)-category encoding. The final synthetic input
always follows the fixed \(17+3+1\) schema. The physical profile is applied to
every informed task.

\subsection{Feature-Profile Provenance}
\label{sec:feature-profile-provenance}

The empirical profiles were built from the usable source corpus, rather than
from the \(608\)-row development subset alone. The composition bank contains
distinct compositions from the \(760\) usable rows, and the environmental
profile contains the \(565\) Fe/Ni--Cr--Mo rows from that same corpus. The
stored numeric imputation means and complete method-category mapping were also
constructed using feature information across the usable corpus.

EPIT values are not stored in these feature profiles or used to calculate
their feature statistics. Target availability is used to identify the usable
task rows. Nevertheless, the profiles expose the synthetic generator to
feature information from final-test rows, including compositions assigned to
final-test groups. Composition separation therefore applies to the labeled
development and evaluation partitions; it does not imply that final-test
compositions were absent from the synthetic feature bank. This makes the prior
dataset-specific and limits a strictly inductive interpretation of the result.

\begin{figure}[htbp]
    \centering
    \projectfigure{images/corrpfn_empirical_feature_distributions.pdf}
    {Observed and generated composition distributions, material-family
    frequencies, and environmental tuple distributions.}
    \caption{Comparison of empirical source features and generated informed
    features. Composition perturbation changes positive element amounts while
    preserving template sparsity; environmental tuples are sampled intact.}
    \label{fig:feature-distributions}
\end{figure}

\subsection{Raw and Standardized Feature Representations}
\label{sec:feature-standardization}

A separate standardized representation is constructed using the population
mean and standard deviation of each column within the generated task:
\begin{equation}
X_{\mathrm{std},ij}=
\begin{cases}
\displaystyle\frac{X_{\mathrm{raw},ij}-\mu_j}{\sigma_j},
&\sigma_j>10^{-6},\\[6pt]
0,&\sigma_j\leq10^{-6}.
\end{cases}
\label{eq:feature-standardization}
\end{equation}
The raw table is preserved and used by the analytical EPIT rule. The
standardized table is used by both the SCM target head and the transformer.

This separation is necessary because the two target mechanisms have different
requirements. The randomly initialized SCM should receive standardized
features so that variables with large numerical units do not dominate only
because of their scale. In contrast, the analytical EPIT rules must receive
the raw physical values because they contain composition thresholds,
chromium--molybdenum terms, logarithmic chloride transformations, temperature
thresholds, pH corrections, and test-method category lookups.

For informed tasks, the general \texttt{Reg2Cls} preprocessing operation is
bypassed. The explicitly constructed \(X_{\mathrm{std}}\) is passed directly
to the transformer. Constant columns become zero but are not deleted, so the
fixed column positions remain valid. A column with one value contains no
within-task variation; in another task where it varies, it is standardized
normally and can influence the target.

Feature permutation is disabled. The model receives anonymous numeric columns
rather than textual feature names, but their positions remain stable. This
permits position-dependent regularities and deliberately narrows the informed
prior to the evaluated schema. Generic tasks retain their separate existing
preprocessing path.

\subsection{SCM Target from the Informed Features}

The standardized empirical features are passed through a direct non-causal
target head:
\begin{equation}
y_{\mathrm{SCM}}=f_{\theta}^{\mathrm{MLP/tree}}(X_{\mathrm{std}}).
\label{eq:scm-target}
\end{equation}
The supplied \(21\) physical columns are treated as external causes. The
target head produces one response per row without generating a replacement
feature table or modifying the supplied physical features.

The MLP uses the existing MLP-SCM sampling of depth, width, activation,
initialization, dropout structure, and noise. Its final layer produces one
target value per row. When the tree branch is selected, the same standardized
features are passed through the existing direct tree-SCM transformation.
These randomly sampled mappings provide synthetic nonlinear and piecewise
relationships; they are not calibrated mechanistic EPIT equations.

\subsection{Calibrated EPIT Target Component}

Independently of the SCM target head, one calibrated analytical rule \(g_k\)
is sampled and evaluated on the raw physical table:
\begin{equation}
y_{\mathrm{EPIT}}=g_k(X_{\mathrm{raw}}).
\label{eq:epit-target}
\end{equation}
All rows in a synthetic task use the same rule family and its saved
coefficient vector. These corrosion-target families are distinct from the Fe
and Ni--Cr--Mo alloy families used by the feature generator.

The seven rules include linear PREN-inspired material resistance, Cr--Mo/W
synergy, chromium thresholds and Mo/W saturation, separate environmental
penalties, coupled environmental breakdown, Fe/Ni-specific thresholds, and a
test-method correction. Their calibration and sampling probabilities are
given in Section~\ref{sec:model-selection}; their equations are collected in
Appendix~\ref{app:target-rules}.

The material and environmental terms encode broad tendencies rather than
universal physical laws. Chloride concentration and, in many systems,
temperature tend to reduce EPIT, but their effects depend on the alloy and
experimental range. The influence of pH is likewise alloy- and range-dependent.
The candidate rules consequently include different environmental structures.
PREN-inspired material terms omit nitrogen because it is absent from the
retained schema.

\subsection{SCM--EPIT Target Mixture}
\label{sec:target-mixture}

Both target components are standardized before mixing, and the combined
target is standardized once more:
\begin{equation}
y_{\mathrm{informed}}
=\zscore\!\left[
 (1-\lambda)\zscore(y_{\mathrm{SCM}})
 +\lambda\zscore(y_{\mathrm{EPIT}})
\right],\qquad0\leq\lambda\leq1.
\label{eq:scm-epit-target}
\end{equation}
Here, \(\zscore\) denotes population standardization within the generated
task, with numerical handling for effectively constant signals. The parameter
\(\lambda\) is implemented as \texttt{informed\_target\_mix\_weight}.

For nondegenerate components, \(\lambda=0\) gives a pure SCM target,
\(\lambda=1\) gives a pure calibrated EPIT-rule target, and intermediate
values combine stochastic SCM diversity with the calibrated corrosion
contribution. Equal mixture coefficients do not necessarily imply equal
fractions of explained target variance, because the components can be
correlated. Final target standardization removes any literal millivolt
interpretation from the synthetic response.

The implementation requires calibrated rule scores in this mode and fails if
they are absent. Unusable synthetic tasks are rejected by the generation
validity checks. The relative weights within an EPIT rule are its calibrated
coefficients; \(\lambda\) controls the complete EPIT component relative to the
SCM component.

\subsection{Generic and Informed Task Mixture}

The hybrid prior samples informed tasks with probability
\[
\rho=\texttt{informed\_prior\_ratio}.
\]
Conceptually, the training-task distribution is
\begin{equation}
\mathcal{D}_{\mathrm{train}}\sim
\begin{cases}
\mathcal{D}_{\mathrm{SCM+EPIT}},&\text{with probability }\rho,\\
\mathcal{D}_{\mathrm{generic}},&\text{with probability }1-\rho.
\end{cases}
\label{eq:hybrid-routing}
\end{equation}
The implementation samples routing and several generator hyperparameters at
the subgroup level. Datasets within a subgroup share these choices while
retaining separately sampled data and noise. This behavior is inherited from
the TabICL training implementation.

For informed tasks, \texttt{mlp\_prob} controls whether the MLP or tree head
generates \(y_{\mathrm{SCM}}\). It has no effect on whether the task is
informed. Generic tasks retain the fixed mixture
\[
P(\mathrm{MLP}\mid\mathrm{generic})=0.7,
\qquad P(\mathrm{tree}\mid\mathrm{generic})=0.3.
\]
Thus, varying the informed MLP/tree mixture does not change the generic
control distribution. Generic tasks receive neither empirical physical
features nor a calibrated EPIT target component.

\subsection{Optional Magpie Representation}

Ten deterministic Magpie-style composition descriptors are available as a
recorded final-training override. When enabled, they are calculated from the
raw composition and appended to the physical table before column-wise
standardization, giving \(31\) model inputs. The SCM target head continues to
use only the \(21\) physical features, and the analytical EPIT rule continues
to use their raw physical representation. The first \(21\) standardized
transformer inputs therefore remain identical to the SCM inputs.

Magpie descriptors are disabled in the base Optuna search and in the reported
Trial~56 model. Their definitions and evaluation requirements are retained in
Appendix~\ref{app:magpie}; they do not contribute to the reported final result.

\section{Target-Rule Calibration and Model Selection}
\label{sec:model-selection}

\subsection{Workflow and Data Boundaries}

The workflow has five stages: composition-aware split construction,
target-rule calibration and direct evaluation, Optuna trial training and
development-fold evaluation, final long training of the selected
configuration, and evaluation of the locked model on the final-test set.
Split construction is described in Section~\ref{sec:composition-split}.

The split manifest records all row assignments, composition components, and
fold identifiers. A separate lock records hashes of the manifest and related
artifacts. Source-data and split hashes are verified before later stages use
them, preventing a newly generated split from being silently combined with
coefficients calibrated on an older split. Inspection files redact final-test
EPIT values. The relevant scripts and artifacts are listed in
Appendix~\ref{app:implementation}.

\subsection{Target-Rule Families}

A common target-rule interface is used for the seven candidates in
Table~\ref{tab:rule-families}. All candidates are evaluated on the same
\(452\) eligible Fe and Ni--Cr--Mo development rows.

\begin{table}[htbp]
\centering
\caption{Candidate analytical EPIT rules used by the informed target sampler.}
\label{tab:rule-families}
\small
\begin{tabularx}{\linewidth}{lL}
\toprule
Rule & Main structure \\
\midrule
Linear PREN & Cr--Mo--W resistance with environment and
material--chloride terms. \\
Cr--Mo/W synergy & Additional Cr--Mo/W cooperation and
temperature--chloride interaction. \\
Threshold and saturation & Smooth chromium threshold and a saturating
Mo/W contribution. \\
Improved environment & Separate logarithmic-chloride,
high-temperature, temperature--chloride, and acidic-pH penalties. \\
Coupled breakdown & Environmental aggressiveness opposed by material
resistance within one breakdown term. \\
Fe/Ni chromium threshold & Separate chromium thresholds for Fe and
Ni--Cr--Mo alloys, with the improved environmental terms. \\
Method-aware & Cr--Mo/W synergy with a test-method correction. \\
\bottomrule
\end{tabularx}
\end{table}

The method-aware calibration uses four broad method families:
potentiodynamic, potentiostatic, scratch, and other. It learns shrunken method
offsets from context targets only. Synthetic generation instead samples
random offsets over the encoded test-method categories; it does not copy the
fitted real-data method directions. The calibrated coefficient controls the
relative contribution of this synthetic method term. Scan rate is excluded
because it is not part of the fixed \(21\)-feature schema.

The historical all-alloy PREN rule and its Fe/Ni-restricted counterpart are
retained only as direct-calibration references. They do not enter the seven-rule
synthetic sampler. Exact rule definitions and the final coefficient vectors
are given in Appendix~\ref{app:target-rules}.

\subsection{Fold-Local Calibration}
\label{sec:rule-calibration}

Target-rule calibration has two separate purposes: estimating rule quality
and determining the coefficient values used during synthetic training.

First, each candidate rule is evaluated using the five fixed development
folds. For fold \(f\), its coefficients and preprocessing state are fitted
using only the eligible context rows. The fitted rule is then applied to the
eligible validation rows, which were not used to fit those coefficients.

The calibration target is the standardized rank of the observed EPIT value.
Let \(B_k\) contain the signed, standardized terms of rule \(k\), let
\(r_y\) be the standardized target ranks, and let \(a_k\) be a predefined
coefficient anchor. The calibration solves
\begin{equation}
\widehat c_k=\arg\min_{c}
\left\{
\frac{1}{n}\|B_kc-r_y\|_2^2+\gamma\|c-a_k\|_2^2
\right\},
\quad \sum_jc_j=1,\quad 0\leq c_j\leq u_{k,j}.
\label{eq:rule-calibration}
\end{equation}
Here, \(\gamma\) controls regularization toward the anchor, and
\(u_{k,j}\) contains the configured upper bounds. Nonnegative coefficients
preserve the signs already assigned to the rule terms. They are dimensionless
relative weights, not fitted electrochemical slopes in millivolts.

For rule \(k\), the direct validation score on fold \(f\) is
\[
r_{k,f}=\operatorname{Spearman}
\left(s_{k,f}(X_{\mathrm{val},f}),y_{\mathrm{val},f}\right),
\]
where \(s_{k,f}\) is the rule fitted on that fold's context portion. The
rule-quality score is
\begin{equation}
s_k=\frac{1}{5}\sum_{f=1}^{5}r_{k,f}.
\label{eq:rule-quality}
\end{equation}
This fold procedure is necessary because fitting and scoring a rule on the
same rows would exaggerate its apparent quality.

After direct evaluation, each rule is fitted once more using all eligible
development rows. For the seven candidates, this means all \(452\) eligible
rows, not the complete \(608\)-row development set. These final-development
coefficients are frozen and reused during both Optuna pretraining and final
long training.

Thus, the five-fold mean Spearman score determines how often a rule is sampled,
while the final-development fit supplies the coefficient values used to
generate synthetic targets. Context-fitted real-data scaling supports direct
rule validation; synthetic rule terms are standardized within each generated
task. The final-development coefficients are transferred to those synthetic
terms rather than used as a direct final-inference predictor.

\subsection{Direct Rule Results and Sampling Probabilities}
\label{sec:rule-probabilities}

Table~\ref{tab:rule-probabilities} reports the development-fold scores and the
resulting sampling probabilities. All seven scores are positive. Rule \(k\)
is sampled with probability
\begin{equation}
p_k=\frac{s_k}{\sum_{\ell=1}^{7}s_\ell}.
\label{eq:rule-probabilities}
\end{equation}
This normalization gives better-performing rules a larger sampling
probability while keeping all seven candidates present. It is a calibrated
task-sampling policy, not Bayesian posterior weighting over corrosion laws.

\begin{table}[htbp]
\centering
\caption{Mean direct development-fold Spearman scores and normalized
synthetic sampling probabilities for the seven candidate target rules.}
\label{tab:rule-probabilities}
\begin{tabular}{lrr}
\toprule
Rule & Mean Spearman & Sampling probability \\
\midrule
Linear PREN & 0.5174 & 0.1342 \\
Cr--Mo/W synergy & 0.5583 & 0.1448 \\
Threshold and saturation & 0.5181 & 0.1343 \\
Improved environment & 0.5850 & 0.1517 \\
Coupled breakdown & 0.5525 & 0.1433 \\
Fe/Ni chromium threshold & 0.5252 & 0.1362 \\
Method-aware & \textbf{0.5999} & \textbf{0.1556} \\
\bottomrule
\end{tabular}
\end{table}

The historical all-alloy PREN baseline obtained a mean fold Spearman
correlation of \(0.4598\). On the matched \(452\)-row Fe/Ni scope, the
historical PREN baseline obtained \(0.5123\). The matched value is the
appropriate reference for the seven candidates. These direct correlations
measure formula alignment, not transformer performance.

The training command supplies all seven scores and all \(29\) calibrated
coefficient values using round-trip-exact floating-point precision. Their
values are therefore identical to those recorded in the calibration artifacts
and pipeline fingerprint; table entries in this report are rounded for
readability.

\begin{figure}[htbp]
    \centering
    \projectfigure{images/corrpfn_rule_calibration.pdf}
    {Five development-fold scores for each analytical rule, their means, and
    the resulting sampling probabilities.}
    \caption{Direct rule calibration on the eligible development rows. Fold
    predictions use context-fitted coefficients; mean fold scores determine
    the synthetic rule-sampling probabilities.}
    \label{fig:rule-calibration}
\end{figure}

\subsection{Optuna Search and Proxy Training}

The final study is identified as
\path{epit_pipeline_optuna_empirical_features_scm_target_v7}.
Its four active search dimensions are shown in
Table~\ref{tab:search-space}.

\begin{table}[htbp]
\centering
\caption{Search space for the final CorrPFN prior. The MLP probability applies
only to informed target heads.}
\label{tab:search-space}
\begin{tabular}{ll}
\toprule
Parameter & Search space \\
\midrule
Informed-task probability \(\rho\) & \(\{0.25,0.50,0.75,1.00\}\) \\
Informed MLP probability \(p_{\mathrm{MLP}}\)
    & \(\{0,0.25,0.50,0.70,0.75,1.00\}\) \\
SCM--EPIT mixture weight \(\lambda\) & \([0,1]\), continuous \\
Composition perturbation \(\tau\) & \([0,0.15]\), continuous \\
\bottomrule
\end{tabular}
\end{table}

The empirical feature profiles, Fe/Ni template-frequency mixture, calibrated
rule scores and coefficients, fixed feature schema, preprocessing policy,
and generic MLP/tree mixture remain fixed. Magpie descriptors are disabled.
Physical feature generation is always active for informed tasks, feature
permutation is disabled, and feature-block coupling is fixed to zero.

The study uses the tree-structured Parzen estimator sampler with ten startup
trials. Each proxy trial is pretrained once from scratch to step \(1{,}000\),
while the learning-rate scheduler retains a \(10{,}000\)-step horizon. The
checkpoint is then evaluated on each of the five saved development folds.
The five validation Spearman correlations are averaged to obtain the trial
objective. MAE, RMSE, \(R^2\), and fold variability are recorded as secondary
diagnostics.

Spearman correlation is the primary metric because the EPIT target covers a
wide numerical range and contains extreme values. Magnitude-based errors can
be strongly influenced by a small number of observations, whereas Spearman
measures whether predicted and observed values have a consistent monotonic
ordering. MAE, RMSE, and \(R^2\) remain necessary because good ranking alone
does not guarantee accurate predictions on the millivolt scale.

Workers can contribute to a shared Optuna journal using one GPU and one
training process per worker. The study records parameters, generated commands,
metrics, seeds, worker configuration, and artifact hashes. The final base
configuration uses NumPy and Torch seeds of \(42\) and one parallel
prior-generation job. Worker and seed settings are part of the recorded
experiment rather than implicit defaults.

\subsection{Interpretation of Development Scores}

The direct rule evaluation is fold-local: each fold's direct rule prediction
is produced using coefficients fitted without that fold's validation targets.

The subsequent Optuna evaluation is intentionally not fully fold-pure. The
sampling probabilities use the direct outcomes from all five development
folds, and the synthetic coefficients are fitted on all eligible development
targets before Optuna begins. Those frozen values are then reused while
every Optuna fold is evaluated. The target rules are not refitted inside each
transformer fold evaluation.

This design keeps the informed-target distribution identical across trials
and excludes outer final-test targets from model selection. It also means
that the Optuna and checkpoint-development scores are model-selection scores
rather than unbiased estimates of generalization. They are reported separately
from the final-test metrics.

\subsection{Final Training and Checkpoint Selection}

Final training selects a completed trial with a finite development objective.
The implementation permits selecting the best completed trial available at
that time even when other trials remain running or waiting; their identifiers
and states are recorded. This is a snapshot decision: later study results
cannot silently change an already frozen final run.

The selected configuration is retrained from scratch for \(10{,}000\)
Stage~1 steps using the same calibrated rule distribution. Transformer weights
are optimized on synthetic tasks. Real development labels influence the
calibrated prior and are supplied as context during evaluation; they are not
used as a conventional supervised transformer training set.

Permanent checkpoints are saved every \(500\) steps. Each candidate is
evaluated on the five frozen development folds using the same inference
configuration. The checkpoint with the highest mean development-fold Spearman
correlation is selected. MAE, RMSE, and \(R^2\) are retained as diagnostics
but do not determine the selection.

A final-model manifest records the selected trial, effective training
configuration, target-rule artifacts, split hashes, checkpoint evaluations,
selected-checkpoint hash, and final inference settings. A separate lock binds
the manifest to that checkpoint. Existing final manifests and locks are not
overwritten. Recorded overrides use a distinct run identity so that they can
be distinguished from an unmodified Optuna winner.

\subsection{Final Inference Configuration}

The selected model receives all \(608\) development rows and their EPIT
targets as context and predicts the \(152\) final-test rows. The locked
configuration uses eight estimators, no feature-order shuffling, only the
numerical view selected by \texttt{norm\_methods=[``none'']}, median
regression output, no regression-uncertainty output, and all available rows.
The normalization option disables additional power-transformed inference
views; it does not remove the estimator's ordinary numeric preprocessing.

The evaluator verifies the model and split hashes and the expected feature
schema before prediction. A completion manifest prevents an existing
successful final evaluation from being silently replaced. Separate comparison
and diagnostic outputs are distinguished from model selection.

The corrosion formulas are not applied directly during final inference.
Their role is to shape the synthetic pretraining distribution. During final
evaluation, the trained TabICL model itself predicts the EPIT values.

\section{Final Results}
\label{sec:results}

\subsection{Selected Prior Configuration}

At the time of final-model freezing, the Optuna study contained \(78\)
trials, of which six were still marked as running. Final training explicitly
selected the best completed trial available at that time, Trial~56. Its mean
five-fold development Spearman correlation was \(0.6903\), with mean MAE
\(211.0\)~mV, RMSE \(298.0\)~mV, and \(R^2=0.4373\).

\begin{table}[htbp]
\centering
\caption{Selected prior configuration for the reported CorrPFN model.}
\label{tab:selected-configuration}
\begin{tabular}{lr}
\toprule
Parameter & Selected value \\
\midrule
Optuna trial & 56 \\
Informed-task probability \(\rho\) & 0.75 \\
Informed MLP probability \(p_{\mathrm{MLP}}\) & 0.50 \\
SCM--EPIT mixture weight \(\lambda\) & 0.2048 \\
Composition perturbation \(\tau\) & 0.1028 \\
Magpie descriptors & disabled \\
Physical feature count & 21 \\
Final training steps & 10,000 \\
Development-selected checkpoint & 2,500 \\
\bottomrule
\end{tabular}
\end{table}

Trial~56 therefore combines \(75\%\) informed tasks with \(25\%\) generic
tasks. Within the informed branch, MLP and tree target heads have equal
probability. The SCM component receives the larger target-mixture coefficient,
while the calibrated EPIT component contributes a smaller explicit domain
bias. These values describe the selected configuration; they are not an
estimate of the proportion of real EPIT variation explained by each mechanism.

\begin{figure}[htbp]
    \centering
    \projectfigure{images/corrpfn_optuna_parallel_coordinates.pdf}
    {The four searched parameters and development Spearman scores of completed
    trials available at final-model freezing, highlighting Trial 56.}
    \caption{Final Optuna search. Each completed trial is represented by its
    prior configuration and mean five-fold development Spearman correlation
    after \(1{,}000\) proxy-training steps.}
    \label{fig:optuna-search}
\end{figure}

\subsection{Development-Based Checkpoint Selection}

Trial~56 was retrained from scratch for \(10{,}000\) Stage~1 steps. Permanent
checkpoints were saved every \(500\) steps and evaluated using the five
frozen development folds. The selected and cryptographically locked checkpoint
was step \(2{,}500\), which obtained mean development-fold Spearman
\(0.6966\), MAE \(191.3\)~mV, RMSE \(280.9\)~mV, and \(R^2=0.4717\).

\begin{table}[htbp]
\centering
\caption{Development metrics for Trial~56 at proxy selection and after final
training. These are selection scores, not final-test results.}
\label{tab:development-results}
\begin{tabular}{lrrrr}
\toprule
Checkpoint & Spearman & MAE (mV) & RMSE (mV) & \(R^2\) \\
\midrule
Proxy run, step 1,000 & 0.6903 & 211.0 & 298.0 & 0.4373 \\
Final run, step 2,500 & 0.6966 & 191.3 & 280.9 & 0.4717 \\
\bottomrule
\end{tabular}
\end{table}

The final training budget does not imply that the last checkpoint is the
preferred model. Development evaluation selected an earlier checkpoint using
the stated metric and without consulting final-test targets.

\begin{figure}[htbp]
    \centering
    \projectfigure{images/corrpfn_development_checkpoints.pdf}
    {Mean development-fold Spearman over the 500-step checkpoints of final
    Trial 56 training, with step 2,500 marked as selected.}
    \caption{Development-only checkpoint selection during final training.
    The reported checkpoint is the one selected by mean development-fold
    Spearman correlation.}
    \label{fig:development-checkpoints}
\end{figure}

\subsection{Comparison Models}

The \emph{generic retrained baseline} uses the same local TabICL architecture
and core training setup, but is pretrained on the unmodified generic
\texttt{mix\_scm} prior with MLP and tree probabilities of \(0.7\) and
\(0.3\). It provides a control for the complete change in pretraining
distribution. It is shown at the matching training step; this comparison
does not establish its best possible checkpoint under independent selection.

The \emph{pretrained TabICL V2} checkpoint provides an external pretrained
model reference. It is not an architecture-matched control because CorrPFN
uses the local TabICL V1 implementation and may lack later model and training
changes included in V2. \emph{CatBoost} provides a conventional supervised
tabular-regression reference. Pretrained V2 and CatBoost are
checkpoint-independent references, evaluated on the same development/final
partition. The comparison uses the \(21\)-feature representation without
Magpie augmentation.

\subsection{Held-Out EPIT Results}

The locked step-\(2{,}500\) CorrPFN checkpoint was evaluated using all
\(608\) development rows as context and the \(152\) composition-separated
final-test rows as queries. Table~\ref{tab:final-results} reports the results.

\begin{table}[htbp]
\centering
\caption{Performance on the \(152\)-row final EPIT test partition. The
CorrPFN checkpoint was selected on development folds; the generic local
baseline is shown at the matching training step.}
\label{tab:final-results}
\begin{tabular}{lrrrr}
\toprule
Model & Spearman & MAE (mV) & RMSE (mV) & \(R^2\) \\
\midrule
CorrPFN, step 2,500
    & \textbf{0.7681} & \textbf{186.8} & \textbf{268.7} & \textbf{0.5294} \\
Pretrained TabICL V2 & 0.6011 & 206.8 & 343.8 & 0.2294 \\
CatBoost & 0.6090 & 231.7 & 323.1 & 0.3194 \\
Generic baseline, step 2,500 & 0.2829 & 326.4 & 548.8 & -0.9630 \\
\bottomrule
\end{tabular}
\end{table}

Relative to pretrained TabICL V2, CorrPFN reduced MAE by \(9.7\%\), reduced
RMSE by \(21.9\%\), and increased Spearman correlation by \(0.1670\).
It also outperformed the reported CatBoost and generic local references
across all four metrics. The result combines improved ordering with improved
millivolt-scale errors, even though configuration and checkpoint selection
were based on Spearman correlation.

Final-test Spearman was higher than the development-fold mean, and MAE and
RMSE were lower. These evaluations use different query rows and different
context sizes: the final model receives \(608\) context rows rather than
\(486\)--\(487\). Their difference should not be interpreted as an isolated
effect of additional context or as evidence of statistical significance.

\begin{figure}[htbp]
    \centering
    \projectfigure{images/corrpfn_final_predictions.pdf}
    {Observed versus predicted EPIT and residuals for the locked step-2,500
    model on the 152 final-test rows.}
    \caption{Final-test predictions of the development-selected CorrPFN
    checkpoint. Predictions and residuals are shown on the original
    millivolt scale.}
    \label{fig:final-predictions}
\end{figure}

\subsection{Post-Evaluation Checkpoint Diagnostics}

All-checkpoint results on the final-test set were inspected afterward for
diagnostic purposes. Step \(2{,}000\) produced the strongest observed
final-test metrics, but it was not selected independently of those targets.
It must therefore not replace the development-selected step-\(2{,}500\)
checkpoint as the reported model. Subsequent decisions informed by these
diagnostics would require a new independent evaluation to support a fresh
confirmatory claim.

\section{Discussion and Limitations}
\label{sec:discussion}

\subsection{What the Final Construction Achieves}

The final construction implements a clear division of responsibilities: the
empirical generator supplies physically structured features, the SCM supplies
stochastic nonlinear task diversity, the calibrated EPIT rule supplies
corrosion-specific target structure, and \(\lambda\) determines their
relative weighting. The separate routing parameter \(\rho\) preserves
generic synthetic tasks within the same pretraining workflow.

Generating physical features before the target avoids a mismatch in which a
synthetic target depends on latent variables whose information is lost during
later composition remapping. Using identical standardized physical columns
for the SCM and transformer keeps the target-generating representation aligned
with the observed inputs. Preserving raw values for the EPIT rules allows
their thresholds and transformations to retain their intended interpretation.

The final model improves all four reported metrics relative to the evaluated
references. This supports the usefulness of the complete informed
pretraining configuration for the present EPIT task. It does not separately
establish the contribution of empirical features, target-rule calibration,
generic-task mixing, or each individual rule. Those claims would require
matched ablations of the final system under the same evaluation protocol.

\subsection{Relevant Lessons from Earlier Development}

Direct synthetic-to-real surrogate alignment helped identify weaknesses in
earlier target rules, but it did not reliably rank prior configurations by
their downstream transformer performance. The final workflow therefore uses
direct calibration to define the analytical rule distribution and actual
transformer development evaluation to select the remaining prior parameters.

Earlier composition transformations could produce unrealistic element
combinations or discard information used by an already generated SCM target.
The final generator addresses these structural issues through empirical
templates and feature-first target construction. These observations explain
the design, but the historical comparisons used different configurations and
evaluation settings and are not treated as controlled ablations of the final
model. The historical experiment record is retained separately in
\texttt{main2.tex}.

\subsection{Statistical and Evaluation Limits}

Only one final training seed was evaluated. Differences between models are
therefore point estimates for the recorded runs, without repeated-training
evidence that would establish their stability. The Optuna proxy also uses
\(1{,}000\) steps while final training runs to \(10{,}000\); development
checkpoint selection limits the reliance on the last checkpoint but does not
prove that the proxy identifies the best possible long-training prior.

The final-test labels are reserved from calibration and model selection in
the final workflow. However, the complete \(760\)-row corpus influenced
historical prior development before the revised split was created. Moreover,
the empirical profiles and their stored feature statistics include feature
information across that corpus. Some final-test compositions can therefore
appear as unlabeled templates during synthetic pretraining. The result
demonstrates prediction across separated labeled composition groups under
this dataset-adapted prior, not evaluation on compositions entirely absent
from every source of pretraining information.

The composition grouping does not establish generalization to independent
publications, laboratories, or alloy databases. Reference overlap across
partitions is allowed, and all observations come from the same source
dataset. A stronger test would use a separate corpus or newly collected
measurements, with feature profiles and model choices frozen beforehand.

\subsection{Physical and Representation Limits}

The empirical generator is not a thermodynamic, phase-stability,
heat-treatment, microstructure, or passivation model. Its composition
diversity is bounded by the observed template bank. In particular, the
Ni--Cr--Mo family contains only \(17\) distinct templates and is much less
diverse than the Fe-alloy family. Other alloy classes occur in the evaluated
task but are not directly represented in the informed composition generator
or the seven-rule calibration scope.

Sampling composition and environment independently removes the source
dataset's specific pairing between alloys and experimental campaigns. It
also assumes that recombining these observed components provides useful
training tasks; it does not establish the feasibility or experimental
relevance of every generated combination. The synthetic features and the
calibrated rules are physically informed rather than a complete physical
description of pitting.

The test-method code has no physical ordering, yet the SCM and transformer
receive it as a single numeric feature. The method-aware rule treats it as a
category, but its synthetic offsets are randomly sampled rather than copied
from calibrated real protocol effects. Finally, fixed feature order makes the
model specific to the documented schema, and synthetic target standardization
removes any direct millivolt interpretation of the generating equations.

The current evidence supports CorrPFN as a working dataset-specific EPIT
predictor. Repeated training seeds, matched final-system ablations, and an
external evaluation with development-only feature profiles would determine
how much of the observed improvement persists beyond this setting.

\appendix

\section{Analytical EPIT Rules and Calibrated Coefficients}
\label{app:target-rules}

\subsection{Notation and Material Scores}

The following definitions describe the weighted analytical rules used in the
final informed branch. Let \(\mathcal{Z}\) denote population standardization
of a signal across a generated task, with a denominator floor of \(10^{-6}\)
and exactly constant signals mapped to zero, and let
\(\sigma(u)=(1+e^{-u})^{-1}\). Material inputs use raw
weight-percent values; \(T\) is in degrees Celsius and \(C_{\mathrm{Cl}}\)
is chloride concentration in molar units. Nonnegative composition inputs and
a positive floor of \(10^{-12}\) for the chloride logarithm are used.

Define the combined molybdenum/tungsten concentration
\[
Q=\mathrm{Mo}+0.55\mathrm{W}.
\]
The four material scores are
\begin{align}
R_{\mathrm{linear}}
    &=\mathrm{Cr}+3.25Q,\label{eq:rule-linear}\\
R_{\mathrm{synergy}}
    &=\mathrm{Cr}+3.25Q+\sqrt{\mathrm{Cr}\,Q},\label{eq:rule-synergy}\\
R_{\mathrm{threshold}}
    &=\sigma\!\left(\frac{\mathrm{Cr}-12}{2}\right)
      +\log(1+Q),\label{eq:rule-threshold}\\
R_{\mathrm{Fe/Ni}}
    &=\sigma\!\left(\frac{\mathrm{Cr}-t_i}{2}\right)
      +\log(1+Q),\label{eq:rule-fe-ni}
\end{align}
where \(t_i=12\) for Fe-alloy rows and \(t_i=15\) for Ni--Cr--Mo rows.
Thresholds, slopes, and the implicit unit concentration in the logarithm are
defined for inputs expressed in weight percent. These expressions are
heuristic response scores rather than dimensionally calibrated physical
equations.

The final empirical generator supplies alloy-family metadata to determine
\(t_i\). Direct calibration uses the recorded material class. The broader
rule interface has a fallback based on whether Ni exceeds Fe, but the final
empirical branch uses the supplied family metadata.

For the material score \(R_k\) associated with rule \(k\), define
\begin{equation}
m_k=\mathcal{Z}(R_k),\qquad
P_k=\mathcal{Z}(\tanh m_k),\qquad
I_k=\mathcal{Z}\!\left(\sigma(-m_k)C\right),
\label{eq:rule-material-signals}
\end{equation}
where \(C\) is the chloride drive defined below. \(P_k\) is the material
passivity term and \(I_k\) is the susceptibility--chloride term.

\subsection{Environmental and Method Terms}

Let
\begin{align}
L&=\mathcal{Z}\!\left(\log_{10}\max(C_{\mathrm{Cl}},10^{-12})\right),\\
C&=\sigma(L),\\
E&=\sigma\!\left(\mathcal{Z}\!\left[
 0.075\mathcal{Z}(T)+0.925L+0.115\mathcal{Z}(|\mathrm{pH}-7.25|)
\right]\right),\\
H&=\sigma\!\left(\frac{T-50}{10}\right),\qquad
A=\sigma(6.5-\mathrm{pH}),\\
J_0&=\sigma(\mathcal{Z}(T))C,\qquad J_1=HC.
\label{eq:rule-environment-signals}
\end{align}
Here, \(E\) is the combined chloride-dominant environment term; \(H\) is
a high-temperature signal; \(A\) is a one-sided acidic-pH signal; and
\(J_0\) and \(J_1\) are two temperature--chloride interactions. The
absolute pH departure in \(E\) and the one-sided acidic term in \(A\)
represent different candidate assumptions, not a universal pH law.

The coupled-breakdown rule uses
\begin{equation}
D_k=\sigma\!\left[
0.65L+0.25\mathcal{Z}(H)+0.10\mathcal{Z}(J_1)-0.75m_k
\right].
\label{eq:rule-coupled}
\end{equation}
This single term combines environmental aggressiveness with material
protection, replacing separate environment and material--chloride penalties.

For the method-aware synthetic rule, one standard-normal offset \(o_j\) is
sampled for each of the \(52\) encoded test-method categories. A row with
category \(a_i\) receives the standardized method term
\begin{equation}
M_i=\mathcal{Z}\!\left[\tanh\!\left(\mathcal{Z}(o_{a_i})\right)\right].
\label{eq:rule-method}
\end{equation}
Both standardizations act on the row-wise signal across the generated task.
The offsets have no imposed physical ordering and are resampled for each
synthetic mechanism.

\subsection{Seven Rule Vectors}

Each rule is a weighted sum of signed, standardized terms:
\begin{equation}
g_k(X)=b_k(X)c_k,
\qquad c_{k,j}\geq0,\qquad\sum_jc_{k,j}=1.
\label{eq:rule-vector}
\end{equation}
The material score and ordered term vector for each rule are as follows:
\begin{enumerate}
    \item \textbf{Linear PREN} (\texttt{pren\_linear}) uses
    \(R_{\mathrm{linear}}\) and
    \[
    b_k=(P_k,-\mathcal{Z}(E),-I_k).
    \]
    \item \textbf{Cr--Mo/W synergy} (\texttt{cr\_mow\_synergy}) uses
    \(R_{\mathrm{synergy}}\) and
    \[
    b_k=(P_k,-\mathcal{Z}(E),-I_k,-\mathcal{Z}(J_0)).
    \]
    \item \textbf{Threshold and saturation}
    (\texttt{threshold\_saturation}) uses \(R_{\mathrm{threshold}}\) and
    the same ordered term vector as the synergy rule.
    \item \textbf{Improved environment}
    (\texttt{improved\_environment}) uses \(R_{\mathrm{linear}}\) and
    \[
    b_k=(P_k,-\mathcal{Z}(C),-\mathcal{Z}(H),
                 -\mathcal{Z}(J_1),-\mathcal{Z}(A)).
    \]
    \item \textbf{Coupled breakdown} (\texttt{coupled\_breakdown}) uses
    \(R_{\mathrm{linear}}\) and
    \[
    b_k=(P_k,-\mathcal{Z}(D_k),-\mathcal{Z}(A)).
    \]
    \item \textbf{Fe/Ni chromium threshold}
    (\texttt{fe\_ni\_cr\_threshold}) uses \(R_{\mathrm{Fe/Ni}}\) and
    the same ordered term vector as the improved-environment rule.
    \item \textbf{Method-aware} (\texttt{method\_aware}) uses
    \(R_{\mathrm{synergy}}\) and
    \[
    b_k=(P_k,-\mathcal{Z}(E),-I_k,-\mathcal{Z}(J_0),M).
    \]
\end{enumerate}

For these weighted rule families, the numerical constants in the equations
above are fixed. The calibrated outer coefficient vectors are also fixed
across training tasks. Rule selection, generated feature tables, synthetic
method offsets, and SCM mechanisms provide task-to-task variation.

For direct real-data rule evaluation, standardization states are fitted on
the relevant context rows and reused on validation rows. For synthetic
generation, the terms are standardized within the generated task. The
method-aware direct calibration additionally replaces random synthetic
offsets with context-fitted offsets for the four broad real method families.
Those offsets are derived from standardized-rank residuals relative to the
anchored synergy rule and shrunk by the factor \(n_j/(n_j+20)\), where
\(n_j\) is the number of context rows in the method family. The fitted offsets
are then centered by their context-row frequencies. Unseen method families
receive zero offset before the fitted method-term scaling is applied.

\subsection{Final-Development Coefficient Vectors}

Table~\ref{tab:rule-coefficients} gives the coefficient vectors stored for
the final run, in the term order defined above. These are the fitted relative
weights used during synthetic generation, not the separate rule-sampling
probabilities in Table~\ref{tab:rule-probabilities}.

\begin{table}[htbp]
\centering
\caption{Frozen final-development rule coefficients, rounded to four decimal
places. Dashes denote terms absent from a rule. Full-precision values are
stored in the calibrated artifacts and final-model manifest.}
\label{tab:rule-coefficients}
\small
\begin{tabular}{lrrrrr}
\toprule
Rule & \(c_1\) & \(c_2\) & \(c_3\) & \(c_4\) & \(c_5\) \\
\midrule
Linear PREN & 0.5560 & 0.4013 & 0.0427 & -- & -- \\
Cr--Mo/W synergy & 0.4331 & 0.0710 & 0.2072 & 0.2887 & -- \\
Threshold and saturation & 0.4186 & 0.1139 & 0.1960 & 0.2715 & -- \\
Improved environment & 0.5394 & 0.2414 & 0.0798 & 0.1393 & 0.0000 \\
Coupled breakdown & 0.4265 & 0.5062 & 0.0673 & -- & -- \\
Fe/Ni chromium threshold & 0.5233 & 0.2721 & 0.1058 & 0.0988 & 0.0000 \\
Method-aware & 0.3051 & 0.0000 & 0.3124 & 0.2099 & 0.1727 \\
\bottomrule
\end{tabular}
\end{table}

Some candidate terms receive zero fitted weight. In particular, the acidic-pH
term is inactive in the final improved-environment and Fe/Ni-threshold fits,
and the combined environment term is inactive in the method-aware fit.
Their rule definitions remain available, but these particular coefficients
should not be described as active contributions in the selected prior.

\section{Optional Magpie-Style Descriptors}
\label{app:magpie}

\subsection{Descriptor Definitions}

The additional features are deterministic composition descriptors rather than
learned material embeddings. They contain no directly trained corrosion
knowledge and do not replace the original element concentrations. The original
\(17\) material features remain model inputs in their existing fixed order.

Let \(c_{ij}\) be the weight-percentage value of element \(j\) in synthetic
or real alloy \(i\), and let \(M_j\) be its atomic weight. A copy of the
material vector is converted to atomic fractions according to
\begin{equation}
a_{ij}=\frac{c_{ij}/M_j}{\sum_e c_{ie}/M_e}.
\label{eq:magpie-atomic-fractions}
\end{equation}
This conversion is used only for descriptor calculation. It does not modify
the weight-percentage features supplied in the original material columns.
The descriptors are calculated from the \(17\) observed elements, so their
atomic fractions refer to the retained composition subset.

For an elemental property \(q_j\), its composition-weighted mean and average
absolute deviation are
\begin{equation}
\mu_q^{(i)}=\sum_j a_{ij}q_j,
\qquad
\delta_q^{(i)}=\sum_j a_{ij}|q_j-\mu_q^{(i)}|.
\label{eq:magpie-mean-deviation}
\end{equation}
Property ranges are calculated as
\begin{equation}
R_q^{(i)}=\max_{j:a_{ij}\geq0.01}q_j
         -\min_{j:a_{ij}\geq0.01}q_j.
\label{eq:magpie-range}
\end{equation}
The \(1\%\) atomic-fraction threshold is used only for range descriptors.
Weighted means and deviations continue to use all elements, and the threshold
does not alter the generated compositions.

\begin{table}[htbp]
\centering
\caption{The optional ten-feature Magpie-style descriptor bundle.}
\label{tab:magpie-descriptors}
\begin{tabularx}{\linewidth}{lL}
\toprule
Elemental property & Calculated statistics \\
\midrule
Electronegativity & Mean, range, average absolute deviation \\
Covalent radius & Mean, average absolute deviation \\
Melting temperature & Mean, range \\
\(d\)-valence electrons & Mean \\
Total valence electrons & Mean \\
Unfilled valence states & Mean \\
\bottomrule
\end{tabularx}
\end{table}

The elemental property and atomic-weight values were extracted once from
Matminer~0.10.0 and stored in a fixed, versioned lookup table identified as
\path{epit_magpie_v1_matminer_0_10_0}. Matminer is therefore used as the
reference source for the elemental values but is not a runtime dependency
during synthetic training. The descriptors are calculated using vectorized
Torch operations during synthetic generation and the corresponding NumPy
implementation during real-data evaluation.

\subsection{Training and Evaluation Integration}

When enabled, the ten descriptors are appended to the raw physical table:
\[
X_{\mathrm{raw}}^{31}
=[X_{\mathrm{raw}}^{21},X_{\mathrm{Magpie}}^{10}].
\]
The complete \(31\)-column table is then standardized column by column before
it is passed to the transformer. Because this operation is column-wise, its
first \(21\) columns are numerically identical to the standardized
\(21\)-column table used by the SCM:
\[
X_{\mathrm{transformer},1:21}^{31}=X_{\mathrm{std}}^{21}.
\]

The descriptors are calculated after the SCM--EPIT target has been constructed
from the \(21\) physical columns. They receive no independent SCM generation,
within-block coupling, random marginal transformation, categorical conversion,
or independent noise. Their relationship with the target arises through their
deterministic dependence on the same composition used by the existing target
mechanisms.

Generic hybrid tasks receive no Magpie descriptors. Their tensors can be
zero-padded to the common batching width when necessary; padded positions are
not interpreted as generated material descriptors.

Descriptor-enabled checkpoints require the same augmentation of the real
EPIT table. For every evaluation split, missing values in the \(17\)
material columns are imputed using means fitted on the training portion of
that split. The same training-fitted means are then applied to the test
portion. The imputed weight percentages are converted to atomic fractions,
and the same ten descriptors are calculated in the same order. The original
\(21\) columns remain first.

The implementation was checked against Matminer reference values for a
multielement alloy, with agreement to numerical precision. The Torch and NumPy
implementations also produce matching descriptor values. Further checks
verify that descriptor augmentation leaves the original physical columns
unchanged and that training and evaluation use the same schema. These
implementation checks establish numerical consistency; the final reported
model does not use the descriptor option.

\section{Implementation and Reproducibility}
\label{app:implementation}

\subsection{Active Configuration}

The final informed path is activated through
\path{pitting_composition_mode=empirical_features_scm_target} within the
\texttt{hybrid\_scm} prior. The implementation names corresponding to the
four searched parameters are
\begin{center}
\begin{tabular}{ll}
\toprule
Symbol & Implementation parameter \\
\midrule
\(\rho\) & \texttt{informed\_prior\_ratio} \\
\(p_{\mathrm{MLP}}\) & \texttt{mlp\_prob}, mapped to
\texttt{informed\_mix\_probs} \\
\(\lambda\) & \texttt{informed\_target\_mix\_weight} \\
\(\tau\) & \texttt{pitting\_composition\_perturb\_strength} \\
\bottomrule
\end{tabular}
\end{center}

The fixed physical feature count is \(21\), with allocation \(17,3,1\).
The informed branch uses the fixed EPIT schema, complete empirical physical
features, no feature permutation, no shared block coupling, and no Dirichlet
redraw. The generic MLP/tree probabilities remain \(0.7,0.3\). The saved
calibrated rule scores and coefficient vectors are required inputs.

\subsection{Source Modules and Workflow Entry Points}

The principal source files are listed below using repository-relative paths.
Historical generator branches coexist in the repository but are not required
to understand the method described in this report.
\begin{itemize}
    \item \path{src/tabicl/prior/epit_composition_profile.py} loads and
    samples the versioned composition templates.
    \item \path{src/tabicl/prior/epit_feature_profile.py} and
    \path{src/tabicl/prior/epit_feature_generator.py} load the environmental
    profile and construct complete physical feature rows.
    \item \path{src/tabicl/prior/dataset.py} integrates empirical generation,
    SCM target heads, weighted EPIT rules, target mixing, and task routing.
    \item \path{src/tabicl/prior/mlp_scm.py} and
    \path{src/tabicl/prior/tree_scm.py} implement the stochastic target heads.
    \item \path{scripts/epit_pipeline/target_rules.py} and
    \path{scripts/epit_pipeline/calibrate_target_rules.py} implement direct
    rule calibration and its saved artifacts.
    \item \path{scripts/epit_pipeline/prepare_splits.py},
    \path{scripts/epit_pipeline/run_optuna.py},
    \path{scripts/epit_pipeline/train_final.py}, and
    \path{scripts/epit_pipeline/evaluate_final.py} implement the split,
    search, final-training, and final-evaluation stages.
    \item \path{scripts/epit_pipeline/evaluate_optuna_folds.py} evaluates
    transformer development folds. The empirical-SCM-target search and
    final-training Slurm launchers select the final informed mode.
    \item \path{src/tabicl/prior/magpie_features.py} implements the optional
    descriptor extension.
\end{itemize}

\subsection{Frozen Artifacts and Result Provenance}

The split artifacts are stored under
\path{corrosion_datasets/analysis/epit_pipeline/splits_v2/}.
They include \texttt{split\_manifest.json}, \texttt{split\_lock.json},
\texttt{split\_assignments.csv}, and \texttt{split\_report.html}. They record
row assignments and component membership, bind the source and split through
hashes, and provide inspection outputs with final-test targets redacted.

The reported model is bound by the manifest and lock in
\begin{quote}
\small\raggedright
\path{corrosion_datasets/analysis/epit_pipeline/final_v1/epit_pipeline_optuna_empirical_features_scm_target_v7/}
\end{quote}
The \texttt{final\_model\_manifest.json} records Trial~56, its effective
parameters, the frozen rule distribution, inference configuration, and
development-checkpoint evaluations. The corresponding
\texttt{final\_model\_lock.json} binds that manifest to the selected
step-\(2{,}500\) checkpoint. The checkpoint resides under
\begin{quote}
\small\raggedright
\path{checkpoints/epit_pipeline_final_v1/final_epit_pipeline_optuna_empirical_features_scm_target_v7_trial_0056/step-2500.ckpt}
\end{quote}

The recorded development-checkpoint results are stored alongside the final
manifest. The reported final-test comparison can be traced to the
step-\(2{,}500\) rows in
\begin{quote}
\small\raggedright
\path{corrosion_datasets/analysis/epit_pipeline/trial56_final_test_all_checkpoints_11172/summary.csv}
\end{quote}
This comparison output also includes the subsequent all-checkpoint
diagnostics. Its additional test results are not the checkpoint-selection
criterion; the selected step and its development score are established by
the final-model manifest.

\subsection{Implementation Validation}

Validation of the final construction covered the following invariants:
\begin{itemize}
    \item empirical compositions are nonnegative, preserve absent elements,
    and close in the complete \(24\)-element space before projection;
    \item environmental tuples match stored empirical records, and their
    sampling is independent of composition-family selection;
    \item the analytical EPIT rule receives raw physical values;
    \item the SCM and transformer receive identical standardized physical
    columns, including zero-valued constant columns at fixed positions;
    \item the informed branch bypasses generic \texttt{Reg2Cls}
    preprocessing while generic tasks retain their established path;
    \item the target equation agrees with an independent calculation for all
    seven rules at \(\lambda\in\{0,0.37,1\}\);
    \item end-to-end generation works with both target heads and with
    \(21\) physical or \(31\) descriptor-augmented model inputs;
    \item accepted tasks have finite features and nonconstant targets, with
    final targets approximately centered and population-standardized;
    \item saved rule scores and coefficients are reconstructed exactly in
    generated training commands;
    \item repeated generation with identical NumPy, Torch, and Python seeds
    reproduces features, targets, and selected rule families; and
    \item split, rule, study, and model hashes remain consistent throughout
    calibration, Optuna reconstruction, checkpoint selection, and evaluation.
\end{itemize}

The recorded implementation validation also included Torch/NumPy rule parity,
reduced end-to-end training, and repository regression checks. These checks
support numerical consistency and reproducibility of the documented code
paths. They do not establish physical validity of every generated task or
replace the statistical limitations of the final experiment.

\begin{thebibliography}{9}
\raggedright
\bibitem{nyby2021electrochemical}
Nyby et al.
\textit{Electrochemical metrics for corrosion resistant alloys}.
Scientific Data (2021).
\href{https://doi.org/10.1038/s41597-021-00840-y}{doi:10.1038/s41597-021-00840-y}.
Associated dataset:
\href{https://figshare.com/articles/dataset/Electrochemical_Metrics_for_Corrosion_Resistant_Alloys/13038257}{Electrochemical Metrics for Corrosion Resistant Alloys, Figshare}.
\end{thebibliography}

\end{document}
